% Nastavení dokumentu
\documentclass[a4paper,12pt]{report}
\usepackage[utf8]{inputenc}
\usepackage{hyperref}
\usepackage{graphicx}
\usepackage[table]{xcolor}
\usepackage{tabularray}
\usepackage{geometry}

% Vlastní příkazy
\newcommand{\chp}[1]{\chapter*{#1}\addcontentsline{toc}{chapter}{#1}}
\newcommand{\schp}[1]{\section*{#1}\addcontentsline{toc}{section}{#1}}

% Nastavení Titulní strana: Nadpis, autor, datum
\title{Laboratorní cvičení PWM}
\author{Vojtěch Vašek}
\date{\today}

\begin{document}

% Tisk titulní strany, obsahu a vložení nové stránky
\maketitle
\newgeometry{top=0cm}\tableofcontents\restoregeometry
\newpage

\newgeometry{top=0cm}
\chp{Zadání a úvaha}
Průměrná hodnota napětí na výstupním pinu, ze kterého budeme odebírat sinusový signál, se bude měnit v závislosti na délce trvání logické 1 a logické 0 během jedné časové periody.

\chp{Použité přístroje}
\begin{enumerate}
	\item[•]{Digitální osciloskop - max input: 30 Vpp, MX output: 10 Vpp, 2× 12 MHz PC-Scope → používaný vstup 1}
	\item[•]{Programátor - propojení desky s počítačem; AVR-ISP}
	\item[•]{deska s ATmega16 - typ krystalu: HC49/S QM - 14.7456 MHz, pin A7, potenciometr}
	\item[•]{Počítač - Windows 10 21H2; Kate; AVR Studio 4; AVR\_ISP\_prog; PCLab2000LT}
\end{enumerate}
\restoregeometry
\schp{Obrázek zapojení filtru}
\begin{center}
\includegraphics[width=11cm]{MIT\_21\_4\_25@2.png}
\end{center}

\newgeometry{top=0cm}
\chp{Postup práce}

Po vytvoření nového projektu v AVR Stuidu jsem si zkopíroval halvičkový soubor základního templatu. V hlavním souboru jsem následně includnul zmíněný hlavičkový soubor a napsal jsem návěší \verb|Main| podprogramu.
Následně v Helixu jsem pomocí datasheetu nakonfiguroval v \verb|Reset| osbluze časovač a napsal jsem další části kódu.

V JavaScriptu jsem si napsal kód generující tabulku s hodnoty vzorků přímo naformátovanou pro použití v ASM kódu.

Následně jsem testoval kód a dolaďoval chyby v kódu.

\schp{Zdrojový kód}
Pro lepší přehlednost tohoto dokumentu jsem se rozhodl nevkládat zdrojový kód přímo zde, ale odkázat na něj pomocí \href{https://github.com/Honeymoon-with-Anxiety/Honeymoon/tree/E4A/MIT/21_3_25}{URL}.

Kód jsem si rozdělil na \verb|templateM16|, kde mám \verb|RESET| obsluhu a konfigurace a \verb|21_3_25|, kde mám \verb|Main| smyčku a tabulku vzorků.

\restoregeometry
\newgeometry{top=0cm}
\chp{Obrázek z osciloskopu}
\includegraphics[width=10cm]{MIT\_19\_4\_25@2.png}

\restoregeometry
\includegraphics[width=11cm]{MIT\_19\_4\_25@3.png}
\restoregeometry

\end{document}
